\chapter{Vitiligo}
\index{Vitiligo}

\section*{Definition and Overview}
Vitiligo is an acquired, chronic, progressive cutaneous disorder characterized by the selective destruction of functioning melanocytes, resulting in well-demarcated depigmented macules and patches of the skin, hair follicles, and mucous membranes. The term is believed to originate from the Latin word \textit{vitium}, meaning ``blemish'' or ``defect'', or possibly from \textit{vituli}, referring to the white patches on the skin of calves. Historically, vitiligo has been documented for thousands of years, with mentions in ancient texts such as the Ebers Papyrus of ancient Egypt and the Rig Veda of ancient India, where it was frequently conflated with leprosy (\textit{Kustha}), leading to severe, unwarranted social stigmatization that persists in many cultures today.

In the paradigm of modern dermatology, vitiligo is recognized not merely as a cosmetic concern, but as a complex, multifactorial, systemic autoinflammatory and autoimmune disease. The progressive loss of epidermal melanocytes halts the synthesis of melanin—the primary pigment responsible for the color of the skin, hair, and eyes. This depigmentation deprives the skin of its natural photoprotective shield against ultraviolet (UV) radiation. The disorder is broadly classified into two main clinical categories: non-segmental vitiligo (NSV), which is the most common form and is typically characterized by symmetrical, bilateral lesions that spread over time, and segmental vitiligo (SV), which is characterized by a unilateral distribution matching a specific dermatome or neural segment and tends to stabilize rapidly. A subset of patients may exhibit mixed vitiligo, where non-segmental patches develop on top of a pre-existing segmental pattern. The significance of vitiligo lies not only in its cutaneous manifestation but also in its profound psychological impact, causing significant anxiety, depression, and social impairment, as well as its strong pathogenic association with other systemic autoimmune diseases.

\section*{Epidemiology}
The global prevalence of vitiligo is estimated to range between 0.5\% and 2.0\% of the general population, although wide variations exist across different geographic regions and ethnic cohorts. While some regional studies suggest higher prevalence rates—up to 4.0\% in certain communities in India or parts of East Asia and Africa—these disparities are often attributed to higher reporting rates and increased visibility of depigmented lesions on darker skin phototypes, rather than a true biological variation.

Vitiligo affects individuals of all races, ethnicities, and skin types. It exhibits no gender predilection, affecting males and females in equal proportions. However, epidemiological studies show that females, especially young women, present for medical evaluation and psychological support more frequently than males, likely due to heightened societal cosmetic standards and the associated social stigma of skin disfigurement.

The onset of vitiligo can occur at any age, from early infancy to late adulthood. However, in approximately 50\% of patients, the disease manifests before the age of 20, and up to 25\% of cases begin before the age of 10. Late-onset vitiligo, presenting after the age of 50, is less common and is more frequently associated with thyroid disease or other systemic autoimmune conditions.

Heredity plays a key role in susceptibility. Approximately 20\% to 30\% of individuals diagnosed with vitiligo report a positive family history of the disorder in a first-degree relative. The concordance rate among monozygotic twins is approximately 23\%, indicating that while genetics are crucial, environmental triggers and epigenetic factors are necessary to initiate the disease. The inheritance of vitiligo does not follow a simple Mendelian pattern; instead, it is a polygenic, multifactorial trait governed by multiple susceptibility loci that interact with environmental stressors.

\section*{Aetiology and Pathophysiology}
The precise aetiology of vitiligo remains a subject of ongoing scientific investigation, but it is widely accepted to be a multifactorial disorder. The convergence theory suggests that various pathological mechanisms—including autoimmune response, oxidative stress, genetic predisposition, neural factors, and cell adhesion defects—work in synergy to cause melanocyte death.

\begin{itemize}
    \item \textbf{Autoimmune Destruction}: The dominant theory is that vitiligo is an autoimmune disease where cytotoxic CD8+ T cells target and destroy epidermal melanocytes. In active vitiligo, there is a prominent infiltration of these T cells at the margins of depigmenting lesions. These cytotoxic cells recognize specific melanocytic antigens, such as glycoprotein 100 (gp100), Melan-A (MART-1), tyrosinase, and tyrosinase-related proteins 1 and 2 (TRP-1 and TRP-2). This autoimmune mechanism is supported by the frequent co-occurrence of vitiligo with other autoimmune diseases. Up to 15\% to 20\% of patients with non-segmental vitiligo have a co-existing autoimmune condition, most commonly Hashimoto's thyroiditis or Graves' disease, followed by type 1 diabetes, rheumatoid arthritis, alopecia areata, pernicious anaemia, systemic lupus erythematosus, and Addison's disease.
    \item \textbf{Oxidative Stress and Epidermal Microenvironment}: Melanocytes from vitiligo patients are highly vulnerable to oxidative stress. There is an abnormal accumulation of reactive oxygen species (ROS), particularly hydrogen peroxide, in the epidermis of affected individuals. This is coupled with a systemic or localized deficiency in antioxidant enzymes, such as catalase, superoxide dismutase, and glutathione peroxidase. The accumulation of ROS damages lipids, proteins, and DNA within the melanocytes, leading to mitochondrial dysfunction, endoplasmic reticulum stress, and cell death. The cellular damage also triggers the release of damage-associated molecular patterns (DAMPs) and heat shock protein 70 (HSP70i), which activate dendritic cells, bridging the innate and adaptive immune responses.
    \item \textbf{Genetic Susceptibility}: Genome-wide association studies (GWAS) have identified over 50 genetic susceptibility loci associated with vitiligo. The majority of these genes encode proteins involved in the regulation of the immune system and immune-cell survival, such as the Major Histocompatibility Complex (MHC) class I and II genes, PTPN22, NLRP1, CTLA4, IL2RA, and LPP. Other identified susceptibility genes are specific to melanocyte biology, such as TYR (encoding tyrosinase), which suggests that primary melanocyte defects may make them more visible or vulnerable to immune clearance.
    \item \textbf{Neural Hypothesis}: This pathway is especially relevant to segmental vitiligo, which presents in a dermatomal distribution. It suggests that abnormal release of neurochemical mediators (such as catecholamines, neuropeptide Y, substance P, and vasoactive intestinal peptide) from cutaneous nerve endings can exert direct cytotoxic effects on nearby melanocytes. Local neurogenic inflammation and vasoconstriction may also increase local oxidative stress, precipitating melanocytic destruction.
    \item \textbf{Melanocytorrhagy (Cell Adhesion Defect)}: This theory describes the mechanical detachment and subsequent loss of melanocytes from the basal layer of the epidermis. Under physical stress or friction, E-cadherin expression (a transmembrane protein essential for cell-to-cell adhesion) is downregulated in the melanocytes of vitiligo patients. This leads to impaired adhesion to basement membrane proteins (like collagen and laminin) and adjacent keratinocytes. The melanocytes detach and migrate upwards, eventually being shed from the epidermis. This mechanism explains the localized phenomenon of Koebnerization.
\end{itemize}

\section*{Clinical Features}
The hallmark clinical presentation of vitiligo is the appearance of acquired, asymptomatic, chalk-white or milk-white macules (less than 5 mm) and patches (greater than 5 mm) on the skin. The lesions are usually well-demarcated with sharp, distinct margins. While the disease is primarily asymptomatic, some patients report mild pruritus (itching) or a burning sensation immediately before the appearance of a new lesion or during active spreading phases, suggesting a localized inflammatory process.

The clinical presentation of vitiligo is highly variable and can be categorized based on morphology and distribution:

\begin{itemize}
    \item \textbf{Non-Segmental Vitiligo (NSV)}: The most common clinical subtype, accounting for over 85\% to 90\% of all cases. It is characterized by symmetrical, bilateral distribution. NSV is further subdivided into:
    \begin{itemize}
        \item \textit{Generalized}: Widespread distribution, with lesions commonly occurring on the face, neck, trunk, and extensor surfaces of limbs (elbows, knees).
        \item \textit{Acrofacial}: Depigmentation limited to the face (especially around the eyes, nose, and mouth) and distal extremities (fingertips, hands, toes, and feet).
        \item \textit{Mucosal}: Involvement of oral, labial, or genital mucosae, which can occur in isolation or as part of generalized vitiligo.
        \item \textit{Universal}: Extreme progression where depigmentation covers more than 80\% to 90\% of the total body surface area.
    \end{itemize}
    \item \textbf{Segmental Vitiligo (SV)}: Accounts for approximately 5\% to 15\% of vitiligo cases, showing a preference for children and young adults. It presents as one or more depigmented patches restricted to one side of the body in a distribution that often follows a dermatome. Unlike NSV, SV typically has an abrupt onset, spreads rapidly within its segment, and then stabilizes, remaining stationary for the rest of the patient's life.
    \item \textbf{Mixed Vitiligo}: An uncommon presentation where a patient with a stabilized segmental vitiligo lesion subsequently develops bilateral, symmetrical non-segmental lesions elsewhere on the body.
    \item \textbf{Poliosis (Leukotrichia)}: Depigmentation of the hair follicles within vitiligo lesions. This can affect scalp hair, eyebrows, eyelashes, and body hair, turning it white. The presence of leukotrichia indicates destruction of the follicular melanocyte stem cell reservoir, representing a poor prognostic factor for repigmentation treatments.
    \item \textbf{Koebner Phenomenon}: The development of vitiligo lesions at sites of cutaneous trauma, friction, pressure, or irritation. Common triggers include physical cuts, abrasions, thermal burns, chronic friction from tight clothing (e.g., bra straps, belts, footwear), and sunburns.
    \item \textbf{Atypical Variations}:
    \begin{itemize}
        \item \textit{Trichrome Vitiligo}: Shows a transitional zone of hypopigmentation between the normal skin and the completely depigmented center, resulting in three distinct shades of color (normal, light brown, white).
        \item \textit{Quadrichrome Vitiligo}: Includes a fourth shade, typically hyperpigmented margins at the border of the lesion.
        \item \textit{Inflammatory Vitiligo}: Characterized by a raised, erythematous (red), itchy border at the margin of the depigmented patch, representing active lymphocytic infiltration.
    \end{itemize}
\end{itemize}

\section*{Differential Diagnosis}
It is critical to distinguish vitiligo from other hypopigmenting and depigmenting skin disorders to ensure correct management and avoid unnecessary treatment.

\begin{itemize}
    \item \textbf{Pityriasis Alba}: A common, benign, self-limiting dermatosis primarily affecting children and young adolescents, often associated with atopic dermatitis. It presents as poorly demarcated, mildly hypopigmented (rather than completely depigmented) patches, typically on the face, neck, and upper arms. The lesions often have a fine, powdery scale and lack the sharp boundaries and chalk-white fluorescence under Wood's lamp seen in vitiligo.
    \item \textbf{Tinea Versicolor (Pityriasis Versicolor)}: A superficial fungal infection caused by the lipophilic yeast \textit{Malassezia furfur}. It presents as multiple small, round, scaling macules on the trunk, neck, and upper arms. The lesions can be hypopigmented, hyperpigmented, or erythematous. It is distinguished from vitiligo by the presence of fine scaling (which becomes more prominent when scratched), a positive potassium hydroxide (KOH) microscopy showing a classic ``spaghetti and meatballs'' pattern of hyphae and spores, and a coppery-orange fluorescence under Wood's lamp.
    \item \textbf{Post-Inflammatory Hypopigmentation}: A localized loss of pigment that occurs secondary to an antecedent inflammatory skin disease, such as psoriasis, atopic dermatitis, lichen planus, discoid lupus, physical trauma, or chemical burns. Unlike vitiligo, the lesions are hypopigmented (retaining some melanocytes and basal pigment) and follow the distribution of the original inflammatory eruption. Clinical history and the gradual recovery of pigment over months to years help confirm the diagnosis.
    \item \textbf{Chemical Depigmentation (Occupational Vitiligo)}: Caused by direct exposure to specific chemical agents, primarily phenolic and catechol derivatives (e.g., monobenzyl ether of hydroquinone, para-tertiary butylphenol) found in industrial detergents, rubber gloves, adhesives, and hair dyes. The clinical lesions are indistinguishable from vitiligo but are initially localized to the site of chemical contact (such as the hands or forearms) before potentially spreading. Diagnosis relies on a thorough occupational and chemical exposure history.
    \item \textbf{Lichen Sclerosus}: A chronic inflammatory dermatosis that mainly affects the anogenital region. It presents as ivory-white or porcelain-white papules and plaques. Unlike the asymptomatic and flat patches of vitiligo, lichen sclerosus lesions are characterized by epidermal atrophy (giving the skin a wrinkled, ``parchment paper'' or ``cigarette paper'' texture), pruritus, pain, purpura, and sometimes follicular plugging.
    \item \textbf{Piebaldism}: A rare, autosomal dominant congenital disorder of melanocyte development caused by mutations in the \textit{KIT} proto-oncogene. It is present at birth and remains stable throughout life. It is characterized by a distinctive white forelock of hair in approximately 90\% of cases, and symmetrical, depigmented patches on the forehead, chest, abdomen, and extremities that typically contain islands of hyperpigmented skin.
    \item \textbf{Leprosy (Hansen's Disease)}: An infectious disease caused by \textit{Mycobacterium leprae} that remains endemic in several developing countries. Tuberculoid leprosy presents as hypopigmented macules or plaques. It is differentiated from vitiligo by a loss of cutaneous sensation (thermal, tactile, and pain) within the lesions, hair loss (alopecia) over the patches, anhidrosis (absence of sweating), and thickening of peripheral nerves.
    \item \textbf{Tuberous Sclerosis Complex}: An autosomal dominant neurocutaneous syndrome characterized by the development of benign tumors in multiple organs. The earliest skin manifestations are ``ash-leaf'' spots—hypopigmented, lanceolate macules present at birth or early infancy. They are hypopigmented rather than depigmented and are accompanied by other signs like angiofibromas, shagreen patches, seizures, and developmental delay.
\end{itemize}

\section*{When to Seek Medical Care}
Although vitiligo is not an acute, life-threatening condition, prompt consultation with a general practitioner or dermatologist is strongly recommended upon the first appearance of unexplained white patches on the skin, hair, or mucous membranes. Early clinical assessment is crucial to establish an accurate diagnosis, rule out mimics, screen for associated systemic autoimmune disorders, and initiate treatment during the active phase of the disease, when interventions are most effective in halting progression and promoting repigmentation.

Patients should seek immediate medical attention if they experience:
\begin{itemize}
    \item Rapidly spreading depigmentation covering large areas of the body within a short timeframe, which may require systemic immunosuppressive therapy to stabilize.
    \item Signs of severe cutaneous inflammation, intense itching, pain, or blistering on depigmented patches, particularly after sun exposure, indicating severe sunburn.
    \item Development of symptoms suggestive of underlying endocrine or systemic autoimmune dysfunction. These symptoms include chronic fatigue, unexpected weight gain or loss, cold or heat intolerance, palpitations, muscle weakness, joint pain, hair loss, or dry skin (indicative of thyroid dysfunction).
    \item Severe psychological distress, depression, clinical anxiety, or social withdrawal stemming from the cosmetic impact of the disease.
\end{itemize}

In rare instances, patients with widespread autoimmune disease may develop acute, life-threatening complications related to other autoimmune endocrine glands, such as an adrenal crisis (Addisonian crisis) secondary to autoimmune adrenalitis. Symptoms of an adrenal crisis include severe vomiting, abdominal pain, profound dizziness or fainting (syncope) due to severe hypotension, confusion, and profound weakness. This is a medical emergency. In such cases, contact emergency services immediately:
\begin{itemize}
    \item 	extbf{Local crisis support}: Contact local emergency services
    \item \textbf{United States and Canada}: Call 911
    \item \textbf{United Kingdom}: Call 999 or 112
    \item \textbf{Europe}: Call 112
\end{itemize}

\section*{Diagnosis and Investigations}
The diagnosis of vitiligo is primarily clinical, based on a comprehensive medical history and a detailed physical examination of the skin, hair, and mucous membranes. Laboratory and diagnostic tests are primarily utilized to rule out other conditions and screen for comorbid autoimmune disorders.

\begin{itemize}
    \item \textbf{Clinical History and Physical Examination}: The physician evaluates the pattern of depigmentation, age of onset, rate of progression, presence of the Koebner phenomenon, family history of vitiligo or other autoimmune diseases, and history of chemical exposures. Physical examination involves evaluating the distribution and margin characteristics of the lesions.
    \item \textbf{Wood's Lamp Examination}: A Wood's lamp emits long-wave ultraviolet A (UVA) light (wavelength of 365 nm). In a completely darkened room, vitiligo lesions shine with a bright, chalky, blue-white fluorescence. This examination is essential for identifying lesions in individuals with very fair skin, mapping the extent of the disease, distinguishing depigmentation (complete absence of pigment in vitiligo) from hypopigmentation (partial pigment loss in pityriasis alba or post-inflammatory states), and detecting subclinical lesions that are not yet visible to the naked eye.
    \item \textbf{Thyroid Function and Autoantibody Screening}: Due to the strong association between non-segmental vitiligo and autoimmune thyroid disease, screening is highly recommended. Initial laboratory tests should include serum Thyroid-Stimulating Hormone (TSH) levels, free Thyroxine (T4), and free Triiodothyronine (T3). Testing for thyroid autoantibodies, specifically Anti-Thyroid Peroxidase (TPO) antibodies and Anti-Thyroglobulin (TG) antibodies, is critical, as antibody positivity can precede clinical thyroid dysfunction by years.
    \item \textbf{Skin Biopsy and Histopathology}: A punch biopsy of a lesion is rarely required but may be performed in atypical cases. Histopathological evaluation of an established vitiligo lesion reveals a complete absence of epidermal melanocytes and melanin pigment in the basal layer, which can be confirmed using special stains (such as Fontana-Masson) or immunohistochemistry for melanocyte-specific markers (such as Melan-A, S100, or SOX10). The active border of a spreading lesion may show a sparse-to-moderate inflammatory infiltrate consisting mainly of CD8+ and CD4+ T lymphocytes at the dermo-epidermal junction.
    \item \textbf{Additional Autoimmune Screening}: Depending on clinical suspicion (such as microcytic anaemia, paresthesias, or polyuria), additional testing may include a Complete Blood Count (CBC) and Vitamin B12 levels (to screen for Pernicious Anaemia), fasting plasma glucose or Glycated Haemoglobin (HbA1c) (to screen for Type 1 Diabetes), and an Antinuclear Antibody (ANA) screen.
\end{itemize}

\section*{Management}
The management of vitiligo is multi-faceted and must be tailored to the patient's disease subtype, activity level, distribution, extent of involvement, skin phototype, and psychosocial needs. The goals of therapy are to arrest the active autoimmune destruction of melanocytes, stabilize existing lesions, and promote the migration and proliferation of melanocytes from reservoirs (such as hair follicles) to achieve repigmentation.

\begin{itemize}
    \item \textbf{Topical Pharmacotherapy}:
    \begin{itemize}
        \item \textit{Topical Corticosteroids (TCS)}: High- or medium-potency topical corticosteroids (e.g., clobetasol propionate 0.05\% or mometasone furoate 0.1\% cream) are first-line therapies for localized vitiligo, particularly during active phases. They work by suppressing the localized immune response. Application is typically limited to once daily for a maximum of 3 months, or used in an intermittent schedule (e.g., 5 days on, 2 days off) to minimize local adverse effects such as skin atrophy, striae, telangiectasia, and acneiform eruptions.
        \item \textit{Topical Calcineurin Inhibitors (TCI)}: Tacrolimus (0.03\% or 0.1\% ointment) and pimecrolimus (1\% cream) are immunomodulators that inhibit T-cell activation without causing skin atrophy. They are preferred for sensitive areas such as the face, neck, eyelids, and intertriginous regions (axillae, groin). They are applied twice daily, often in combination with phototherapy or topical corticosteroids.
        \item \textit{Topical Janus Kinase (JAK) Inhibitors}: Ruxolitinib 1.5\% cream is a novel, targeted therapy approved for non-segmental vitiligo. It blocks JAK1 and JAK2 signaling, interrupting the interferon-gamma pathway that drives CD8+ T-cell-mediated destruction of melanocytes. It is applied twice daily and has shown high efficacy, especially on facial lesions.
    \end{itemize}
    \item \textbf{Phototherapy}:
    \begin{itemize}
        \item \textit{Narrowband Ultraviolet B (NB-UVB)}: Administered at a wavelength of 311 nm, NB-UVB is the gold standard for generalized, active non-segmental vitiligo. It has a dual mechanism of action: it suppresses local cutaneous T-cell-mediated immunity and stimulates the proliferation and migration of melanocyte stem cells from the hair follicle bulge. Treatments are typically conducted 2 to 3 times per week, with a minimum of 6 months required to assess efficacy. Maximum therapeutic benefit is often seen after 12 to 18 months of continuous therapy.
        \item \textit{Targeted Phototherapy (Excimer Laser/Lamp)}: The 308 nm excimer laser or monochromatic excimer lamp delivers intense, targeted UVB radiation to localized vitiligo patches, sparing surrounding unaffected skin from hyperpigmentation and UV exposure. It is highly effective for localized lesions on the face and neck.
        \item \textit{Psoralen plus Ultraviolet A (PUVA)}: Involves the administration of a photosensitizing agent (psoralen, either orally or topically as a bath/cream) followed by exposure to UVA light. While effective, PUVA has largely been replaced by NB-UVB due to its higher risk of phototoxic reactions, nausea, headache, and long-term risk of non-melanoma skin cancers.
    \end{itemize}
    \item \textbf{Systemic Therapy (Stabilization of Rapidly Progressing Disease)}:
    \begin{itemize}
        \item \textit{Oral Corticosteroid Mini-Pulses}: For patients with rapidly progressing, generalized vitiligo, low-dose oral corticosteroids (such as dexamethasone 2.5 mg to 5.0 mg or methylprednisolone) are administered on two consecutive days per week (usually weekends) for 3 to 6 months. This pulsatile regimen is highly effective at arresting active disease progression while minimizing the systemic side effects of chronic daily steroid use (e.g., weight gain, osteoporosis, hypertension, hyperglycemia).
        \item \textit{Systemic Immunosuppressants}: Other systemic therapies, including methotrexate, cyclosporine, or oral JAK inhibitors (e.g., tofacitinib), may be considered in refractory, rapidly progressive cases under specialist guidance.
    \end{itemize}
    \item \textbf{Surgical Treatments (for Stable Disease)}:
    Surgical interventions are indicated only for patients with stable vitiligo (defined as no new lesions, no expansion of existing lesions, and no Koebner phenomenon for at least 12 to 24 months) that has failed to respond to medical therapies. These procedures aim to transplant autologous melanocytes from a pigmented donor site (usually the thighs or buttocks) to the depigmented recipient areas.
    \begin{itemize}
        \item \textit{Tissue Grafting}: Includes autologous punch grafting (transferring small skin grafts), suction blister epidermal grafting (harvesting the epidermis using suction blisters to avoid scarring the donor site), and split-thickness skin grafting.
        \item \textit{Cellular Grafting}: Involves harvesting a donor skin sample, separating the cells enzymatically, and preparing a suspension of non-cultured epidermal cells (containing melanocytes and keratinocytes) that is grafted onto a dermabraded recipient site. This technique allows a small donor site to cover a recipient area up to ten times its size.
    \end{itemize}
    \item \textbf{Depigmentation Therapy}:
    For patients with extensive, universal vitiligo (covering more than 80\% of the body surface area) or those with refractory facial lesions who prefer a uniform appearance, permanent depigmentation of the remaining islands of normal skin may be the most viable option. This is achieved using topical application of monobenzyl ether of hydroquinone (monobenzone) in concentrations of 20\% to 40\% applied twice daily for several months. This chemical permanently destroys remaining functional melanocytes. Patients must be counselled on the permanent nature of this therapy and the lifelong requirement for strict sun protection.
\end{itemize}

\section*{Complications}
The complications of vitiligo extend beyond cutaneous depigmentation, encompassing physical, psychological, and systemic risks.

\begin{itemize}
    \item \textbf{Sunburn and Phototoxicity}: Melanin acts as a natural physical filter against ultraviolet radiation. The absolute loss of melanin in vitiligo lesions leaves the skin highly vulnerable to acute actinic damage, leading to severe sunburns, redness, pain, and blistering upon minimal sun exposure. Despite this extreme vulnerability, multiple epidemiological studies have shown that patients with vitiligo do not have a higher incidence of melanoma or non-melanoma skin cancers (such as basal cell carcinoma or squamous cell carcinoma) than the general population. This is hypothesized to be due to the upregulation of tumor suppressor proteins (such as p53) and increased levels of systemic and localized immune surveillance in the skin.
    \item \textbf{Psychosocial and Psychiatric Morbidity}: The psychological impact of vitiligo is profound and frequently underestimated. The visible nature of the disease, particularly in patients with darker skin tones, often leads to significant social stigma, discrimination, teasing, and isolation. Patients have a significantly elevated risk of developing psychiatric comorbidities, including major depressive disorder, generalized anxiety disorder, social phobia, low self-esteem, body dysmorphic disorder, and suicidal ideation. The impact on quality of life is comparable to other chronic diseases like psoriasis, eczema, and type 1 diabetes.
    \item \textbf{Ocular and Auditory Abnormalities}: Because melanocytes are also located in the uveal tract of the eye (iris, ciliary body, and choroid) and the stria vascularis of the cochlea in the inner ear, the systemic autoimmune process in vitiligo can affect these sites. Complications include subclinical uveitis, pigmentary migration in the retina, choroidal atrophy, and mild-to-moderate sensorineural hearing loss. Formal ophthalmological and audiometric evaluations are indicated if symptoms arise.
    \item \textbf{Treatment-Induced Complications}: Prolonged use of high-potency topical corticosteroids can cause local skin atrophy, striae, telangiectasia, hypertrichosis, and systemic absorption leading to hypothalamic-pituitary-adrenal axis suppression. Phototherapy carries risks of acute redness (sunburn-like reactions), itching, dry skin (exfoliation), premature skin aging (photoaging), and potential long-term risks of carcinogenesis (particularly with PUVA). Surgical procedures carry risks of localized infection, graft failure, scarring, hyperpigmentation of donor sites, and an uneven texture at recipient sites.
\end{itemize}

\section*{Prognosis and Outcomes}
Vitiligo is a chronic, lifelong skin condition characterized by an unpredictable clinical course. Complete spontaneous and permanent repigmentation is rare, occurring in less than 1\% to 2\% of cases, and is usually limited to small, localized, sun-exposed areas. For the vast majority of patients, the disease is progressive, with periods of stability punctuated by unpredictable flares of active depigmentation.

The prognosis and response to therapy are heavily influenced by the disease subtype and the anatomical location of the lesions:
\begin{itemize}
    \item Segmental vitiligo (SV) carries a favorable prognosis for stabilization. The depigmentation typically develops rapidly over 3 to 12 months, stays confined to the affected segment, and then ceases to progress. However, SV is highly resistant to medical therapies because the follicular melanocyte stem cell reservoir in the segment is often completely depleted early in the disease course. Consequently, surgical grafting is the treatment of choice once stability is confirmed.
    \item Non-segmental vitiligo (NSV) is characterized by a lifelong risk of progressive spreading. Prognostic indicators of poor response to medical treatment and a higher likelihood of relapse include:
    \begin{itemize}
        \item Early age of disease onset (childhood).
        \item Presence of mucosal involvement (lips, genitals).
        \item Extensive body surface area involvement at presentation.
        \item Presence of poliosis (white hair), which indicates a loss of follicular melanocyte reservoirs.
        \item Lesions located on the hands, feet, fingers, toes, and lips (acral and mucosal sites). These areas naturally contain fewer hair follicles and have a poorer blood supply, making them highly resistant to repigmentation therapies.
    \end{itemize}
\end{itemize}
Repigmentation, when it occurs, typically starts in a ``perifollicular'' pattern, presenting as tiny brown macules around hair follicles within the white patch, which gradually expand and coalesce to cover the depigmented area. A ``diffuse'' or ``marginal'' repigmentation pattern can also occur, spreading inward from the borders of the lesion. While current treatments can achieve significant repigmentation (especially on the face and trunk), they do not cure the underlying genetic and immunological predisposition, and recurrence of depigmentation in previously treated areas occurs in approximately 40\% of patients within 2 to 5 years after discontinuing therapy. Maintenance therapy (such as twice-weekly application of topical tacrolimus) is often recommended to prevent relapses.

\section*{Prevention and Self-Care}
Because the precise cause of vitiligo remains multifactorial, there are no definitive measures to prevent the initial onset of the disease. However, once diagnosed, strict self-care and lifestyle modifications are essential to stabilize the condition, prevent flares, protect vulnerable skin, and manage the psychological impact.

\begin{itemize}
    \item \textbf{Rigorous Sun Protection}: Daily application of a broad-spectrum, water-resistant sunscreen with a Sun Protection Factor (SPF) of 50 or higher is critical. Sunscreen should protect against both UVA and UVB radiation and be applied to all exposed skin, particularly depigmented patches, 20 minutes before going outdoors and reapplied every two hours. Wearing sun-protective clothing (UPF-rated long-sleeved shirts, wide-brimmed hats, and UV-blocking sunglasses) and avoiding outdoor activities during peak UV hours (10:00 AM to 4:00 PM) are crucial. Sunscreen also prevents sunburn, which can trigger the Koebner phenomenon, and minimizes tanning of surrounding normal skin, reducing the visual contrast between pigmented and depigmented areas.
    \item \textbf{Avoidance of Cutaneous Trauma}: Patients must minimize physical trauma, friction, and chemical irritation to the skin to prevent the development of new lesions at injury sites (Koebner phenomenon). Actions include avoiding tight-fitting clothing, restrictive belts, and footwear that rubs against the skin; avoiding aggressive scratching, rubbing, or scrubbing with loofahs; and avoiding cosmetic procedures that damage the epidermis, such as waxing, microdermabrasion, chemical peels, and tattooing.
    \item \textbf{Cosmetic Camouflage}: Cosmetic camouflage techniques can significantly improve quality of life and self-esteem. Options include:
    \begin{itemize}
        \item \textit{Cover-up Makeup}: High-pigment, waterproof camouflage creams and foundations specifically formulated to match the patient's natural skin tone and cover white patches.
        \item \textit{Self-Tanning Products}: Creams or sprays containing dihydroxyacetone (DHA). DHA reacts chemically with keratinocytes in the stratum corneum to produce a temporary brown color that does not wash off with water and lasts for several days, gradually fading as the skin exfoliates.
        \item \textit{Micropigmentation}: Cosmetic tattooing can be used to restore color to stable, small lesions in areas where medical repigmentation is difficult, such as the lips or fingertips. However, it carries a risk of triggering the Koebner phenomenon if the disease is active, and the pigment color may not match changing natural skin tones over time.
    \end{itemize}
    \item \textbf{Stress Management}: Psychological and emotional stress is a well-documented trigger for the onset and exacerbation of vitiligo. Incorporating stress-reduction techniques—such as mindfulness-based cognitive therapy, meditation, yoga, regular physical exercise, and adequate sleep—can help modulate the immune system and reduce disease activity.
    \item \textbf{Nutritional Support}: Eating a balanced diet rich in natural antioxidants (such as vitamins C, E, and beta-carotene, found in fresh fruits, leafy vegetables, nuts, and seeds) may help combat the systemic oxidative stress implicated in melanocyte destruction. Screening for and correcting vitamin D deficiencies is also recommended, as vitamin D plays a role in melanocyte maturation and immune regulation.
    \item \textbf{Psychological Support and Counseling}: Engaging with professional mental health services, such as cognitive behavioral therapy (CBT), and joining patient support groups can help individuals cope with the emotional challenges, anxiety, depression, and social stigma associated with vitiligo.
\end{itemize}

\section*{Sources and Further Reading}
\begin{itemize}
    \item American Academy of Dermatology (AAD). Vitiligo: Overview, Diagnosis, and Treatment. \textit{https://www.aad.org/public/diseases/color-problems/vitiligo}. Accessed August 2026.
    \item National Institute of Arthritis and Musculoskeletal and Skin Diseases (NIAMS). Vitiligo Health Information. \textit{https://www.niams.nih.gov/health-topics/vitiligo}. Accessed August 2026.
    \item British Association of Dermatologists (BAD). Vitiligo Patient Information Leaflet. \textit{https://www.bad.org.uk/patient-information-leaflets/vitiligo}. Accessed August 2026.
    \item Vitiligo Support International (VSI). Patient Support, Advocacy, and Clinical Trials. \textit{https://www.vitiligosupport.org}. Accessed August 2026.
    \item Mayo Clinic. Vitiligo: Symptoms, Causes, Diagnosis, and Treatment. \textit{https://www.mayoclinic.org/diseases-conditions/vitiligo/symptoms-causes/syc-20355912}. Accessed August 2026.
    \item World Health Organization (WHO). Skin Diseases and Autoimmune Disorders. \textit{https://www.who.int}. Accessed August 2026.
\end{itemize}